\section{A product estimate across a fast bridge}
\label{bridge:section}

The second estimate compares observations near opposite ends of a long
interval on which the transfer matrix grows almost at its maximum rate.
The comparison comes from two transverse families of contracting graphs
in the middle of the interval. All measures in this section are the
normalized area measure, expressed in the coordinates
$(q_i,q_{i-1})$ when convenient, whenever they define probabilities
or expectations. Label lengths and slice integrals use ordinary
one-dimensional Lebesgue measure in their displayed coordinates.

For distinct integers $d,e$, write
\begin{equation}
 \mathfrak h_{d,e}
   =1-\frac{\log_M\norm{T_{d,e}}}{\abs{d-e}}.
 \label{bridge:shortfalls}
\end{equation}
These quantities belong to $[0,1]$. An observation with bound $(R,L,B)$
means a function of at most $R$ such quantities whose defining function
on $[0,1]^R$ has absolute value at most $B$ and Lipschitz constant at most
$L$ for the maximum norm. Allowing fewer than $R$ arguments is harmless.
The defining functions and the integer indices may vary from one
application to another, provided these bounds are preserved.
Equivalently, these bounds may be imposed only on the realized range
$\mathcal R\subseteq[0,1]^R$ of the arguments. Indeed, a function $f$
on that range extends by
$u\mapsto\inf_{v\in\mathcal R}(f(v)+L\norm{u-v}_\infty)$;
clipping this extension to $[-B,B]$ preserves its values on
$\mathcal R$ and its Lipschitz bound.

\begin{theorem}[Product estimate across a fast bridge]
\label{bridge:theorem}
There are absolute constants $c_b\in(0,1)$, $\kappa>0$, and $\sigma>0$
with the following property. Fix observation bounds $(R,L,B)$ and a
threshold $h>0$. For sufficiently large $M$ and sufficiently large
$n=b-a$, let $W,H$ be observations with these bounds, with $W\geq0$.
All indices occurring in one observation are within $\kappa n$ of $a$,
and all indices occurring in the other are within $\kappa n$ of $b$.
Either assignment of $W,H$ to the two ends is allowed. Then
\begin{equation}
 \E\!\left[
   \one_{\{\log_M\norm{T_{a,b}}\geq c_b n\}}
   W\one_{\{H>h\}}
 \right]
 \leq
 C\bigl(\E W+M^{-\sigma n}\bigr)\PP(H>h/2).
 \label{bridge:product}
\end{equation}
The constants $c_b,\kappa,\sigma$ are independent of the observation
bounds and threshold. The constant $C$ and the lower size requirements
for $M,n$ may depend on $(R,L,B)$ and $h$, but not on the indices,
the defining functions, or the bridge location.
\end{theorem}

We first construct the graphs needed in the proof. Finite-time
contracting directions and curves have been developed in stable-manifold
constructions under weak hyperbolicity hypotheses
\cite{HollandLuzzatto2006}. Here the graphs are
level sets of a distant orbit coordinate, obtained from finite
Dirichlet problems; we prove their extent and transverse distortion
directly. We use $c=0.98$ throughout the construction. A direction is
understood modulo sign.

\begin{lemma}[Contracting graphs and their distortion]
\label{bridge:graphs}
For sufficiently large $M$, suppose $B_+\geq2$ and, at a fixed state with central coordinates
$(x,y)=(q_0,q_{-1})$,
\begin{equation}
 D_l:=\norm{T_{0,l}}\geq M^{cl}
 \qquad(1\leq l\leq B_+).
 \label{bridge:prefix-growth}
\end{equation}
For at least one choice $j\in\{B_+-1,B_+\}$ there is a smooth graph
through this state, parametrized by $q_{-1}=y+r$, $\abs r\leq1/4$,
on which the lifted coordinate $q_j$ is fixed. On this graph,
\begin{equation}
 \abs{q_p(r)-q_p(0)}+
 \abs{\frac{\dd q_p(r)}{\dd r}}
 \leq C M^{-0.8(p+1)}
 \qquad(0\leq p<j).
 \label{bridge:graph-contraction}
\end{equation}
The corresponding endpoint pair at time $B_+$ changes by at most
$C M^{-0.8(B_+-1)}$. The partial derivative
$\partial q_j/\partial q_0$, with $q_{-1}$ held independent, is
nonzero throughout the graph, and the absolute values of this partial
derivative at any two points of the graph have a ratio between two
absolute positive constants.

For a fixed choice of $j$, one may require the closed eligibility
condition
\begin{equation}
 \abs{\partial q_j/\partial q_0}\geq D_{B_+}/3
 \quad\text{at the initial point}.
 \label{bridge:row-choice}
\end{equation}
The graphs depend continuously on eligible initial points. There is
an identical backward statement, with almost horizontal graphs in
$(x,y)$, obtained by using $(q_{-1},q_0)$ as the ordered initial pair.
\end{lemma}

\begin{proof}
Let $s_l$ be a most contracted unit input direction for $T_{0,l}$.
Since a single inverse step has norm at most $M$,
\[
 \norm{T_{0,l}s_{l+1}}\leq\frac{M}{D_{l+1}}.
\]
The projection onto the expanded input direction for $T_{0,l}$
therefore gives
\begin{equation}
 \abs{s_{l+1}\wedge s_l}\leq\frac{M}{D_lD_{l+1}}.
 \label{bridge:direction-step}
\end{equation}
The absolute sine between two unoriented lines satisfies the triangle
inequality. Summing \cref{bridge:direction-step}, and using
\cref{bridge:prefix-growth}, yields
\[
 \abs{s_{B_+}\wedge s_l}
 \leq C M^{1-c-2cl}
 \qquad(1\leq l\leq B_+).
\]
Resolve $s_{B_+}$ into the singular input directions for $T_{0,l}$.
Using $D_l\leq M^l$ as well as its lower bound gives
\begin{equation}
 \norm{T_{0,l}s_{B_+}}
 \leq M^{-cl}+C M^{l+1-c-2cl}
 \leq C M^{-0.94l}.
 \label{bridge:stable-prefix}
\end{equation}
The last inequality follows from
$1-c+(1-2c)l=0.02-0.96l\leq-0.94l$ for $l\geq1$.

The first component of an input vector becomes the second component
after one step. Thus \cref{bridge:stable-prefix} at $l=1$ shows that
$s_{B_+}$ is nearly vertical. Let $t$ solve the variational recurrence
with $t_{-1}=0$, $t_0=1$. Horizontal input has a projection bounded
away from zero onto the expanded input direction for $T_{0,B_+}$.
For large $M$, at least one of the two coordinates $t_{B_+}$ and
$t_{B_+-1}$ consequently has magnitude at least $D_{B_+}/3$.
Choose the corresponding $j$. This proves that the two closed choices
in \cref{bridge:row-choice} cover the states under consideration.

The row of $T_{0,B_+}$ selecting $q_j$ has norm at least
$\abs{t_j}\geq D_{B_+}/3$, and its value on $s_{B_+}$ has magnitude
at most $D_{B_+}^{-1}$. Its kernel direction therefore differs from
$s_{B_+}$ by $O(D_{B_+}^{-2})$. Normalize that direction so that its
second component is 1, and let $U$ denote its variational solution.
Then $U_{-1}=1$, $U_j=0$, and, for $0\leq p<j$,
\begin{equation}
 \abs{U_p}\leq C M^{-0.9(p+1)}.
 \label{bridge:linear-dirichlet}
\end{equation}
Indeed, apply \cref{bridge:stable-prefix} at $l=p+1$ to its main
component. The error component is bounded by
\[
 C M^{p+1}D_{B_+}^{-2}
 \leq C M^{-0.96(p+1)},
\]
because $p+1\leq B_+$. Normalization changes these bounds only by an
absolute factor.

We now solve the nonlinear boundary problem. Let $\xi_p$ be the
change of the lifted position $q_p$, with boundary data
$\xi_{-1}=r$ and $\xi_j=0$. The tridiagonal matrix on
$0,\ldots,j-1$ with diagonal $v(q_p)$ and adjacent entries $-1$
has determinant $t_j\ne0$. Its inverse is
\begin{equation}
 G_{lp}=t_{\min(l,p)}U_{\max(l,p)}.
 \label{bridge:green}
\end{equation}
For completeness, the proposed column solves the homogeneous
recurrence except at its source. At that source the output is 1,
because the Wronskian
$t_{p+1}U_p-t_pU_{p+1}$ equals its boundary value 1. Both boundary
values of the proposed column are zero. This verifies
\cref{bridge:green} directly.

Subtract the original orbit recurrence from the new one, and move
the nonlinear Taylor remainder to the right side. With its sign
absorbed into $F_p$, the boundary problem becomes
\begin{equation}
 \xi_l=U_l r+\sum_{p=0}^{j-1}G_{lp}F_p(\xi_p),
 \qquad
 \abs{F_p(u)}\leq C M u^2,
 \quad
 \abs{F'_p(u)}\leq C M\abs u.
 \label{bridge:fixed-point}
\end{equation}
Use the finite-dimensional weighted norm
\[
 \norm{\xi}_*=\max_{0\leq l<j}
          M^{0.8(l+1)}\abs{\xi_l}.
\]
On a ball of fixed radius $R_0$, the Lipschitz bound for the nonlinear
part of \cref{bridge:fixed-point}, before weighting its output at
index $l$, is bounded by a constant depending on $R_0$ times
\[
 \sum_{p\leq l}
 M^{p-0.9(l+1)+1-1.6(p+1)}
 +\sum_{p>l}
 M^{l-0.9(p+1)+1-1.6(p+1)}.
\]
Here $\abs{t_p}\leq M^p$ and \cref{bridge:linear-dirichlet} were
used. After multiplying by $M^{0.8(l+1)}$, the exponents are,
respectively,
\[
 -0.1(l+1)-0.6(p+1),
 \qquad
 1.8(l+1)-2.5(p+1).
\]
Their sums tend uniformly to zero as $M$ tends to infinity,
independently of $j$. The linear term $Ur$ has bounded weighted norm
for $\abs r\leq1/4$. Choose $R_0$ larger than twice this bound and
then choose $M$ large. The right side of \cref{bridge:fixed-point}
maps this ball to itself and is a contraction there. Its unique fixed
point gives the desired graph on the whole parameter interval.

The same estimates show that
$I-G\diag(F'_p(\xi_p))$ has a uniformly bounded inverse in this
weighted norm. Multiplication by the original Dirichlet matrix
identifies it with the varied linearized Dirichlet problem.
Consequently that problem is invertible, and the finite-dimensional
implicit function theorem gives smooth dependence on $r$.
Differentiating \cref{bridge:fixed-point} solves with right side $U$,
which has bounded weighted norm. This proves
\cref{bridge:graph-contraction}. The contraction and its fixed point
also depend continuously on the eligible initial data: all the
coefficients are continuous there, and the contraction constant and
ball are uniform. This proves the asserted continuous dependence.

If $j=B_+$, the endpoint pair estimate follows directly from the
fixed boundary value and \cref{bridge:graph-contraction}. If
$j=B_+-1$, the recurrence and the fixed value of $q_{B_+-1}$ give
$\xi_{B_+}=-\xi_{B_+-2}$, proving the same estimate without an extra
factor of $M$.

It remains to control transverse distortion. At the varied orbit let
$t_p(r)$ again have initial values $t_{-1}(r)=0$, $t_0(r)=1$.
Invertibility of the varied Dirichlet problem gives $t_j(r)\ne0$.
The graph derivative $U_p(r)=\dd q_p(r)/\dd r$ has boundary values
1 and 0 at indices $-1,j$ and solves that varied problem. Its Green
diagonal is therefore $t_p(r)U_p(r)$. Differentiating the determinant
of the varied Dirichlet matrix gives the exact identity
\begin{equation}
 \frac{\dd}{\dd r}\ln\abs{t_j(r)}
 =\sum_{p=0}^{j-1}t_p(r)U_p(r)^2\phi''(q_p(r)).
 \label{bridge:determinant-distortion}
\end{equation}
The derivative of a diagonal entry is
$\phi''(q_p(r))U_p(r)$, explaining the second factor of $U_p$.
Since $\abs{t_p(r)}\leq M^p$ for every varied orbit,
\cref{bridge:graph-contraction} bounds the absolute value of
\cref{bridge:determinant-distortion} by
\[
 C\sum_{p\geq0}M^p M^{-1.6(p+1)}M
 =C\sum_{p\geq0}M^{-0.6(p+1)}\leq C.
\]
Integration over $\abs r\leq1/4$ proves the bounded-ratio claim.

For backward evolution set $\widetilde q_p=q_{-1-p}$. The recurrence
has the same form in these variables. At step $l$ its ordered pair
is the swapped original pair at time $-l$, and the initial pair is
$(q_{-1},q_0)$. Swapping coordinates is an isometry, so the backward
growth hypotheses and all norm estimates are identical. The resulting
graphs vary $q_0$ and have small slope in the $q_{-1}$ coordinate.
\end{proof}

\begin{proof}[Proof of \cref{bridge:theorem}]
We first find many usable middle times. On the bridge event choose a
unit vector realizing $\norm{T_{a,b}}$ and follow its norms along the
bridge. The increments of their logarithms to base $M$ are at most 1,
and their total shortfall from 1 is at most $(1-c_b)n$. By
\cref{setup:shortfall-count}, if $c_b$ is sufficiently close to 1,
at least $\eta n$ indices $i$ in the middle third, for an absolute
$\eta>0$, satisfy
\begin{equation}
 \begin{aligned}
 \norm{T_{i,i+l}}&\geq M^{cl}
       &&(1\leq l\leq b-i),\\
 \norm{T_{i,i-l}}&\geq M^{cl}
       &&(1\leq l\leq i-a).
 \end{aligned}
 \label{bridge:middle-eligibility}
\end{equation}
In detail, remove indices admitting a forward interval starting there
or a forward interval ending there whose average log increment is
less than $c$. The counting bound removes at most
$C(1-c_b)n/(1-c)$ indices. At a retained index, the chosen vector
bounds the forward operator norm on either kind of interval from
below. On an interval ending at $i$, use also equality of the forward
and inverse operator norms to obtain the second inequality in
\cref{bridge:middle-eligibility}. For large $n$, the integer rounding
of the middle third does not affect a fixed positive proportion.

It suffices to establish \cref{bridge:product}, with a uniform
constant, when the bridge indicator is replaced by the indicator of
\cref{bridge:middle-eligibility} for one fixed middle index $i$.
Indeed, the bridge indicator is at most $(\eta n)^{-1}$ times the
sum of these middle-index indicators; multiplying by the nonnegative
weight and summing then costs only an absolute constant.

Shift $i$ to time 0. Put $B_+=b-i$ and $B_-=i-a$; both are at least
$n/3-O(1)$. Apply \cref{bridge:graphs} forward and backward. There
are two possible endpoint-coordinate choices in each direction.
We work with one fixed pair of choices and eventually sum over the
four possibilities. Write $Q(x,y)$ for the lifted forward coordinate
held fixed, and $F(x,y)$ for the lifted backward coordinate held
fixed. On forward graphs $Q_x\ne0$; on backward graphs $F_y\ne0$.
For large $M$, all forward slopes $\abs{\dd x/\dd y}$ and backward
slopes $\abs{\dd y/\dd x}$ are less than a common absolute number
$a_0<1/2$. The relevant transverse partial derivatives have uniformly
bounded absolute ratios along each graph.

\paragraph{Transport of endpoint observations.}
Along a forward graph the endpoint pair at time $b$ changes by at
most $CM^{-0.8(B_+-1)}$. A single forward or inverse map is
$M$-Lipschitz on the lifted coordinates. Thus pair differences
throughout the window of radius $\kappa n$ about $b$ are at most
\[
 C M^{\kappa n+O(1)}M^{-0.8(B_+-1)}.
\]
The one-step matrix difference is at most $CM$ times the position
difference. Since forward and inverse product norms agree, the
endpoints of each observation interval may first be put in increasing
order. For a product of length at most $2\kappa n$, the
telescoping product identity therefore bounds its matrix difference
by a length factor times
\[
 C M^{3\kappa n+O(1)}M^{-0.8(B_+-1)}.
\]
Every product norm compared here is at least 1. The elementary
inequalities $\abs{\norm A-\norm B}\leq\norm{A-B}$ and
$\abs{\log u-\log v}\leq\abs{u-v}$ for $u,v\geq1$ consequently
give the same exponential control for the normalized logarithms in
\cref{bridge:shortfalls}. The positive integer interval lengths in
their denominators cause no additional loss.

Choose fixed $\kappa,\sigma>0$ with
\begin{equation}
 3\kappa+\sigma<0.8/3.
 \label{bridge:error-exponents}
\end{equation}
The strict margin absorbs the length factor and the bounded powers
of $M$ when $n$ is large. After increasing the size requirements to
account for the observation bounds, each observation at the forward
end changes by at most
\begin{equation}
 \delta_n=C M^{-\sigma n}
 \label{bridge:observation-error}
\end{equation}
along a forward graph. The same statement holds for an observation
at the backward end along a backward graph. All these estimates are
uniform in the chosen middle index, endpoint-coordinate choices,
and observation indices. Increase $M,n$ further so that
$\delta_n<h/2$.

\Cref{bridge:sweep-figure} shows the two transverse coordinate families
used next. It depicts the change of variables, not the motion of an orbit.
\begin{figure}[tb]
\centering
\begin{tikzpicture}[x=0.82cm,y=0.82cm,>=Latex,font=\small]
  \definecolor{leafblue}{RGB}{28,75,118}
  \definecolor{sweeporange}{RGB}{167,82,36}
  \draw[gray!65] (0,0) rectangle (7,4.5);
  \draw[gray!65,->] (0,0) -- (7.35,0) node[right] {$x$};
  \draw[gray!65,->] (0,0) -- (0,4.85) node[above] {$y$};
  \draw[gray!55,densely dotted] (0,2.25) -- (7,2.25);
  \node[left] at (0,2.25) {$y_c$};
  \draw[gray!60,densely dashed] (5.4,0) -- (5.4,4.5);
  \node[below] at (5.4,0) {$u$};
  \draw[sweeporange,thick]
    (0.2,1.17) .. controls (2.5,1.26) and (4.6,1.43) .. (6.8,1.51);
  \draw[sweeporange,thick]
    (0.2,2.02) .. controls (2.5,2.13) and (4.6,2.40) .. (6.8,2.46);
  \draw[sweeporange,thick]
    (0.2,2.91) .. controls (2.5,3.02) and (4.6,3.16) .. (6.8,3.24);
  \draw[leafblue,very thick]
    (3.0,0.18) .. controls (2.86,1.35) and (3.18,2.69) .. (3.06,4.33);
  \fill[leafblue] (3.07,2.25) circle (1.7pt);
  \node[below right,leafblue] at (3.07,2.25) {$s$};
  \draw[leafblue,thick,{Bar[width=4pt]}-{Bar[width=4pt]}]
    (2.95,1.34) -- (3.13,3.11);
  \node[left,leafblue,fill=white,inner sep=1pt] at (2.75,3.13) {$D_s$};
  \fill[sweeporange] (5.4,2.40) circle (1.7pt);
  \node[right,sweeporange,fill=white,inner sep=1pt]
    at (5.4,2.40) {$\Theta_s(u,y)$};
  \node[above,leafblue] at (3.18,4.32) {forward $Q$-level};
  \node[above,sweeporange] at (5.74,3.25) {backward $F$-levels};
\end{tikzpicture}
\caption{Coordinate schematic for the bridge change of variables, not
an orbit plot. A forward level is labelled at $y=y_c$ by $s$; the
eligible slice $D_s$ is swept across the square along transverse
backward levels to the coordinate $(u,\Theta_s(u,y))$.}
\label{bridge:sweep-figure}
\end{figure}


\paragraph{Forward labels and their area.}
Assume first that $H$ is at the forward end and $W$ at the backward
end. Cover the middle-time torus by the finitely many closed squares
of a grid of side $s_0=1/32$ in a fundamental domain. Overlaps on
boundaries have area zero and cause no problem for an upper bound.
Fix one square $I_x\times I_y$ and choose $y_c\in I_y$. All position
functions in this square and its small neighborhood are computed
from the same real lifted recurrence starting at $(x,y)$.

Every forward graph through an eligible point in the square extends
over all of $I_y$, with a margin, since its parameter range has length
$1/2$ and $s_0<1/4$. Label it by its intercept $s$ at $y=y_c$ and
write it as $(X(s,y),y)$. Two such graphs which meet have the same
value of $Q$ and coincide locally because $Q_x\ne0$. They then
coincide on the whole common interval: the coincidence set is both
open, by implicit uniqueness, and closed, by continuity. In
particular distinct labels have disjoint images at every fixed $y$.

Let $E$ be the labels of forward graphs through forward-eligible
points of this square for which $H\geq h$. We use closed eligibility
conditions, including \cref{bridge:row-choice}. They and $H\geq h$
define a compact subset of the closed square. Continuous dependence
of the graph in \cref{bridge:graphs} makes its intercept a continuous
function of the initial point. Thus $E$ is compact. The graph is
uniquely determined by its intercept, so these constructions define
$X$ unambiguously on $E\times I_y$.

For these labels, every point of the graph over $I_y$ has $H>h/2$ by
\cref{bridge:observation-error}. The identity
\[
 Q(X(s,y),y)=Q(s,y_c)
\]
and implicit differentiation give
\begin{equation}
 X_s(s,y)=\frac{Q_x(s,y_c)}{Q_x(X(s,y),y)}.
 \label{bridge:label-jacobian}
\end{equation}
The derivative in \cref{bridge:label-jacobian} is interpreted through
a local smooth extension when $E$ has empty interior; the existence
of that extension is justified below. The distortion assertion of
\cref{bridge:graphs} bounds its absolute value above and below by
positive constants.

The map $(s,y)\mapsto(X(s,y),y)$ on $E\times I_y$ is injective and
has Jacobian comparable to 1. Its image stays in a fixed bounded
lift: the small slope bounds the displacement from the square by an
absolute constant. Projection of this bounded lift to the torus has
bounded multiplicity. Change of variables therefore yields
\begin{equation}
 \abs E\leq C\PP(H>h/2).
 \label{bridge:label-area}
\end{equation}
The factor $\abs{I_y}=s_0$ has been absorbed into the constant.

\paragraph{A transverse sweep for the weight.}
Fix $s\in E$. Let $D_s\subseteq I_y$ consist of the parameters $y$
for which $z(y)=(X(s,y),y)$ lies in the square and is backward-eligible
in the fixed backward endpoint choice. This is a compact set. Carry
each $z(y)$ along its backward graph to first coordinate $u\in I_x$,
and denote the resulting second coordinate by $\Theta_s(u,y)$.
The carry is defined throughout $I_x$, with a margin.

For each fixed $u$, the map $y\mapsto\Theta_s(u,y)$ is injective on
$D_s$. If two carried points agreed, their backward graphs would
coincide throughout $I_x$ by regular-level uniqueness. A single such
graph cannot meet the chosen forward graph in two distinct points:
two intersections would satisfy both
$\abs{\Delta x}\leq a_0\abs{\Delta y}$ and
$\abs{\Delta y}\leq a_0\abs{\Delta x}$, forcing equality of the
points. This proves the asserted injectivity.

The equation defining the carry is
\[
 F(u,\Theta_s(u,y))=F(X(s,y),y).
\]
Differentiation with respect to the starting parameter gives
\begin{equation}
 \partial_y\Theta_s(u,y)
 =\frac{F_y(z(y))+F_x(z(y))X_y(s,y)}
        {F_y(u,\Theta_s(u,y))}.
 \label{bridge:sweep-jacobian}
\end{equation}
The ratio of the two $F_y$ terms is uniformly bounded above and
below along the backward graph. Moreover $\abs{F_x/F_y}\leq a_0$
at its starting point, since $-F_x/F_y$ is its slope, and
$\abs{X_y}\leq a_0$. Hence the absolute value of
\cref{bridge:sweep-jacobian} is uniformly bounded above and below
by positive constants.

The endpoint observation estimate gives, for all $u\in I_x$ and
$y\in D_s$,
\[
 W(X(s,y),y)\leq W(u,\Theta_s(u,y))+\delta_n.
\]
Apply change of variables to the injective map
$(u,y)\mapsto(u,\Theta_s(u,y))$ on $I_x\times D_s$.
Its image lies in a bounded lift. Since $W\geq0$, its integral over
that image is at most a fixed multiple of $\E W$. Using
\cref{bridge:sweep-jacobian} and then dividing by $\abs{I_x}$ gives
\begin{equation}
 \int_{D_s}W(X(s,y),y)\,\dd y
 \leq C\bigl(\E W+M^{-\sigma n}\bigr).
 \label{bridge:conditional-weight}
\end{equation}

\paragraph{Changes of variables on the eligible subsets.}
Here are details ensuring that the preceding changes of variables
do not require the eligible sets to have interior. Consider first
one forward graph with label $s$. It is a compact regular level arc
over $I_y$ and extends with a margin beyond this interval. The
partial derivative $Q_x$ is nonzero on the arc. The implicit
function theorem in finitely many overlapping charts along the arc
therefore continues it smoothly for intercepts in an open
neighborhood of $s$. The continued level is prescribed by
$Q(s',y_c)$ for each nearby intercept $s'$. Uniqueness in the
overlapping charts makes these local continuations consistent.
Every nearby eligible label follows the same continuation, because
its regular graph passes through that intercept. Thus $X$ on
$E\times I_y$ is locally the restriction of a smooth function of
$(s,y)$, and \cref{bridge:label-jacobian} holds at all required
points.

For the transverse sweep, fix $s$ and a point of $D_s$. Its backward
graph is likewise a compact regular arc, with $F_y\ne0$, extending
with margin over $I_x$. The same continuation argument applies to
the smoothly varying starting points $(X(s,y),y)$ on the fixed
forward graph. It gives a local smooth extension of
$\Theta_s(u,y)$ in both variables, agreeing with all nearby
eligible carries by uniqueness. The maps used in both area
comparisons are consequently restrictions of local smooth
diffeomorphisms, with the displayed nonvanishing Jacobians.

A countable cover by such charts can be split into disjoint Borel
pieces. On each piece ordinary change of variables applies to its
measurable subset; the images of the pieces are disjoint because
the full restricted maps are injective. Summing proves the formulas
on the compact eligible domains. This also establishes the
measurability of all integrands and swept regions used above.
No extension of eligibility to a neighborhood is assumed.

\paragraph{Completion of the area comparison.}
Every point of the fixed square satisfying both eligibility
conditions and $H>h$ belongs to a forward graph with label in $E$.
Change variables from $x$ to this label at each $y$, using
\cref{bridge:label-jacobian}. Drop conditions other than square
membership and backward eligibility, which is allowed because the
weight is nonnegative. The desired weighted integral over this
square and this pair of choices is bounded by
\[
 C\int_E\int_{D_s}W(X(s,y),y)\,\dd y\,\dd s.
\]
Combining \cref{bridge:conditional-weight,bridge:label-area} bounds
this by the right side of \cref{bridge:product}.

If $H$ is at the backward end, begin instead with the labels of the
almost horizontal backward graphs and sweep them in the $x$
direction to bound their label measure by $\PP(H>h/2)$. For each
such graph, transport its eligible points along the almost vertical
forward graphs to estimate $W$. The proof just given applies with
$x,y$ and $Q,F$ interchanged: the graph extents, distortion bounds,
small-slope injectivity argument, and local chart construction are
the same. This proves the estimate for the other assignment of the
two observations as well.

Finally sum over the finitely many squares and the at most four
endpoint-coordinate choices. The constants remain uniform. Average
over the usable middle indices as explained after
\cref{bridge:middle-eligibility}. This establishes
\cref{bridge:product} and completes the proof.
\end{proof}
